\subsection{Aplicaciones: II. Anillos universalmente japoneses.}
\label{IV.2.ss7.7-fr}

\begin{env}[7.7.1]
\label{IV.2.s7.7.1-fr}
Recordemos (\textbf 0, 23.1.1) que se llama anillo \emph{universalmente japonés} a un
anillo $A$ tal que toda $A$-álgebra \emph{íntegra de tipo finito} es un anillo japonés. Es equivalente decir que, para toda
$A$-álgebra íntegra \emph{de tipo finito} $B$, el cierre
entero de $B$ es una $B$-álgebra \emph{finita}.
\end{env}
